\chapter{Analytical Investigation and Design}

\section{Finite Element Model Setup}

\begin{figure}[H]
    \centering
    \includegraphics[width=0.7\textwidth]{cad_isometric.png} 
    \caption{High-fidelity 3D CAD Model of the construction helmet shell, showing exterior ribbing and peak design.}
    \label{fig:cad_model}
\end{figure}

\begin{figure}[H]
    \centering
    \includegraphics[width=0.7\textwidth]{cad_section_liner.png}
    \caption{Section view of the CAD assembly detailing the interface between the exterior polymer shell (yellow), the internal EPS energy-absorbing liner (blue), and the rigid headform (gray).}
    \label{fig:cad_section_liner}
\end{figure}
The simulations were executed using the LS-DYNA explicit time integration 
solver within ANSYS Workbench. As defined by the kinematic methodology in 
Section 3.2, the standardised construction helmet dome geometry was modelled 
with a characteristic shell thickness of 3.0 mm, paired with a standardised 
EPS energy-absorbing liner. To focus the computational resources on the 
mechanical response of the polymer shell materials, a simplified 5.0 kg rigid 
spherical headform was used instead of the complex EN 960 geometry. 
This simplification ensures contact stability at the liner interface and 
provides a consistent kinematic reference point for force transmission 
maintain the exact mass parity required by the EN 397 
impact standard.

\section{Constitutive Material Modelling}
Accurate FEA of transient impacts requires constitutive models capable of capturing both the elastic and post-yield behavior of the polymers. 

\subsection{Thermoplastic Shells (\texttt{*MAT\_024})}
The outer structural shells (Virgin HDPE, ABS, Bamboo-HDPE, and Recycled PC) were modelled using LS-DYNA's \texttt{*MAT\_PIECEWISE\_LINEAR\_PLASTICITY} (\texttt{*MAT\_024})~\cite{lsdyna_keyword}. This constitutes an elasto-plastic material model that evaluates yielding based on the von Mises criterion. The model relies on the input of a defined yield stress ($\sigma_y$) and an arbitrary isotropic hardening curve, which was approximated using the Tangent Modulus ($E_{tan}$). For stress states exceeding the yield threshold, the flow rule allows for permanent plastic deformation while maintaining volume constancy (incompressibility in the plastic domain). This model is highly effective for thin-walled structures undergoing localized yielding, as observed in the HDPE drop tests.

\subsection{Energy-Absorbing Liner (\texttt{*MAT\_063})}
The Expanded Polystyrene (EPS) liner required a fundamentally different mathematical approach. Unlike the incompressible plastic flow of the shell, foams undergo massive permanent volumetric reduction. The liner was modelled using \texttt{*MAT\_CRUSHABLE\_FOAM} (\texttt{*MAT\_063})~\cite{lsdyna_keyword}. This model calculates volume compression separately from shape distortion. 

The compressive behavior is governed by a user-defined stress versus volumetric strain curve. Volumetric strain is defined directly by relative volume compaction ($\varepsilon_v = 1 - V/V_0$). Unloading is treated as purely elastic up to a specified tension cut-off. For this EN 397 simulation, this reflects the physical reality that the EPS liner absorbs the headform's kinetic energy through irreversible cellular collapse during the primary crush phase, but provides minimal structural resistance during the subsequent elastic rebound.

\section{Mesh Convergence Study}

A mesh convergence study was performed on the Virgin HDPE baseline model to ensure that the results were independent of the spatial discretization. Five different element sizes were evaluated (6~mm, 5~mm, 4~mm, 3~mm, and 2~mm), with the peak equivalent stress monitored as the primary convergence metric. The model predominantly uses linear 4-node shell elements with Belytschko-Tsay integration and viscous hourglass control (Type 4) to maintain thermodynamic stability during the high-velocity impact~\cite{belytschko2013, lsdyna2021}.

As shown in Table~\ref{tab:mesh_convergence} and Figure~\ref{fig:convergence}, the maximum equivalent stress stabilized as the element size was reduced. The transition from 3~mm to 2~mm resulted in a change of 5.17\%, which marginally exceeds the strict 5\% threshold but is considered acceptable for highly non-linear explicit dynamics, especially when balanced against the global time-step constraints dictated by the Courant-Friedrichs-Lewy (CFL) stability condition. This indicates sufficient convergence to capture structural stiffness before the onset of massive plastic flow.


\begin{table}[H]
    \centering
    \caption{Mesh convergence study metrics for Virgin HDPE helmet shell (sub-critical impact energy).}
    \label{tab:mesh_convergence}
    \begin{tabularx}{\textwidth}{C C C C}
        \toprule
        \textbf{Mesh Size (mm)} & \textbf{Elements (approx.)} & \textbf{Peak Stress (kg/mm$^2$)} & \textbf{Max Displacement (mm)} \\
        \midrule
        6.0  & 4,800  & 2.40 & 44.02 \\
        5.0  & 6,700  & 2.38 & 40.37 \\
        4.0  & 10,500 & 2.34 & 42.63 \\
        3.0  & 18,900 & 2.01 & 51.56 \\
        2.0  & 42,400 & 2.11 & 44.14 \\
        \bottomrule
    \end{tabularx}
\end{table}

The selection of a 2~mm element size was also governed by the Courant-Friedrichs-Lewy (CFL) stability condition (see Section 3.1.2), which dictates that the solver time-step ($\Delta t$) must not exceed the time taken for a pressure wave to cross the smallest element dimension ($L_{min}$)~\cite{lsdyna2021}. The final model achieved a stable time-step of approx. $0.42~\mu s$. Validation of 0.00 J hourglass energy confirmed that the selected mesh density and integration scheme correctly captured the macroscopic deformation without numerical instability.

\begin{figure}[H]
    \centering
    \includegraphics[width=0.7\textwidth]{mesh_at_crown.png}
    \caption{Finite Element Mesh of the helmet shell highlighting the 2 mm element resolution at the crown impact zone.}
    \label{fig:helmet_mesh}
\end{figure}

\begin{figure}[H]
    \centering
    \includegraphics[width=0.7\textwidth]{mesh_convergence_HDPE.pdf}
    \caption{Mesh convergence study for Virgin HDPE shell showing peak equivalent stress stabilization.}
    \label{fig:convergence}
\end{figure}
\subsection{Mesh Quality Metrics}
To ensure the reliability of the explicit dynamics results, the mesh quality was audited against industry-standard metrics for skewness, orthogonal quality, and aspect ratio~\cite{lsdyna2021}. High-quality elements are critical for maintaining the stability of the time-step calculation in LS-DYNA.

\begin{figure}[H]
    \centering
    \begin{subfigure}[b]{0.48\textwidth}
        \centering
        \includegraphics[width=\textwidth]{skewness.png}
        \caption{Skewness}
        \label{fig:mesh_skewness}
    \end{subfigure}
    \hfill
    \begin{subfigure}[b]{0.48\textwidth}
        \centering
        \includegraphics[width=\textwidth]{orthogonal_quality.png}
        \caption{Orthogonal Quality}
        \label{fig:mesh_ortho}
    \end{subfigure}
    
    \vspace{10pt}
    
    \begin{subfigure}[b]{0.48\textwidth}
        \centering
        \includegraphics[width=\textwidth]{aspect_ratio.png}
        \caption{Aspect Ratio}
        \label{fig:mesh_aspect}
    \end{subfigure}
    \hfill
    \begin{subfigure}[b]{0.48\textwidth}
        \centering
        \includegraphics[width=\textwidth]{element_ratio.png}
        \caption{Element Ratio}
        \label{fig:element_ratio}
    \end{subfigure}
    \caption{Mesh quality distribution for the 2 mm helmet model. The majority of elements maintain an orthogonal quality above 0.8 and a skewness below 0.3, ensuring high numerical accuracy. The element ratio confirms the dominance of high-quality quadrilaterals in the impact zone.}
    \label{fig:mesh_quality}
\end{figure}

\section{Contact Parameters and Interaction}
The interaction between the helmet assembly components was governed by the 
LS-DYNA *Automatic Surface-to-Surface* penalty contact algorithm. A penalty 
stiffness factor of 0.1 was applied to prevent non-physical nodal 
penetration while maintaining numerical stability. Friction was modelled 
using a Coulomb friction coefficient of $\mu = 0.3$ between the polymer shell 
and the rigid headform, and $\mu = 0.5$ between the shell and the impacting 
ceiling, consistent with standard polymer-on-metal contact mechanics in 
crashworthiness simulations~\cite{mishra2021}.

\section{Material Models and Test Matrix}
Four distinct materials were evaluated. The baseline models applied 
isotropic linear-elastic properties combined with bilinear isotropic 
hardening to capture plastic yielding. The material properties used in the 
simulations are summarized in Table~\ref{tab:material_props}.

\begin{table}[H]
    \centering
    \caption[Material properties assigned to the helmet shell FEA models.]{Material properties assigned to the helmet shell FEA models. 
    Values for conventional materials sourced from~\cite{cruz2023} and 
    sustainable alternatives derived from~\cite{chandramohan2015, yadav2024}.}
    \label{tab:material_props}
    \begin{tabularx}{\textwidth}{L C C C C C L}
        \toprule
        \textbf{Material} & \textbf{E (GPa)} & \textbf{Yield (kg/mm$^2$)} & \textbf{Density (kg/m$^3$)} & \textbf{$\nu$} & \textbf{E$_{tan}$ (GPa)} & \textbf{Source} \\
        \midrule
        Virgin HDPE & 1.00 & 2.55 & 950 & 0.42 & 0.10 & \cite{tham2014} \\
        ABS Plastic & 2.10 & 4.08 & 1050 & 0.35 & 0.50 & \cite{rigoni2017} \\
        Bamboo-HDPE & 1.50 & 3.26 & 1020 & 0.38 & 0.30 & \cite{mohanty2010} \\
        Recycled PC & 2.30 & 6.12 & 1200 & 0.37 & 0.60 & \cite{reich2019} \\
        \bottomrule
    \end{tabularx}
\end{table}

The Bamboo-HDPE composite is approximated as an isotropic material. This 
simplification is justified based on the assumption of a random 
short-fibre distribution within the 30\% wt blend, which promotes 
quasi-isotropic macroscopic behaviour (Chandramohan et al., 2015). However, 
potential orthotropy remains a limitation of the study (see 
Section~\ref{sec:limitations}).

\section{Force Extraction and Signal Processing}

To determine the EN 397 regulatory compliance, the transmitted force ($F_t$) was derived from the peak nodal acceleration ($a_n$) of the headform's center of gravity:
\begin{equation}
    F_t = m_{hf} \cdot a_{filt}
\end{equation}
where $m_{hf}$ represents the standardized headform mass (5.0~kg) and $a_{filt}$ is the nodal acceleration signal after the application of the SAE J211 CFC 60 low-pass filter~\cite{sae1995}.


The raw nodal acceleration outputs were processed in Python. The 4-pole 
phaseless Butterworth filter (SAE J211 CFC 60 specification) was applied 
with a cutoff frequency of 100 Hz. The filter was applied bidirectionally 
(forward and reverse pass) to achieve zero phase distortion, consistent 
with the SAE J211 standard requirement for crash test data processing.

\subsection{Signal Processing Refinements}
The raw explicit data initially generated non-physical nodal acceleration spikes. Prior to filtering, the raw deceleration data exhibited severe high-frequency ringing. This was a direct artefact of the explicit penalty contact forces echoing between the rigid 5.0~kg headform and the inner surface of the EPS liner, where stress waves artificially reflect through the rigid elements. Rather than classifying these momentary numerical spikes as structural failures, the 100 Hz cutoff was essential to extract the true macroscopic force required to evaluate the 5.0~kN EN 397 threshold.
Because the raw ANSYS data was sampled at 1000 Hz, a CFC 60 filter (100 Hz 
cutoff) was applied, adhering to the Nyquist--Shannon limit, to 
isolate the true macroscopic deceleration curve.

Deformation tracking was also corrected to measure the \textit{relative} 
kinematic displacement between the headform's centre of gravity and the 
helmet apex, providing the structural crush depth rather than global 
absolute movement.
